\section{Dung dịch, nồng độ, khối lượng riêng và độ tan}

Phần lớn hóa chất trong thực tế được dùng dưới dạng dung dịch. Cần phân biệt rõ \textbf{khối lượng dung dịch}, \textbf{khối lượng chất tan}, \textbf{thể tích dung dịch} và \textbf{số mol chất tan}.

\subsection{Nồng độ phần trăm và nồng độ mol}
\begin{ghinho}
\[
C\%=\frac{m_{\text{chất tan}}}{m_{\text{dung dịch}}}\times100\%,
\qquad
C_M=\frac{n_{\text{chất tan}}}{V_{\text{dung dịch}}(\si{\liter})}.
\]
Nếu biết khối lượng riêng \(D\):
\[
m_{\text{dung dịch}}=DV.
\]
\end{ghinho}

\begin{vidu}
Cần pha \SI{2.0}{\kilo\gram} dung dịch \ce{NaCl} \(0.9\%\) theo khối lượng. Tính khối lượng \ce{NaCl} và nước cần dùng.
\end{vidu}
\begin{loigiai}
\[
m_{\ce{NaCl}}=2.0\times0.009=\SI{18}{\gram}.
\]
Khối lượng nước là \(\SI{2.0}{\kilo\gram}-\SI{0.018}{\kilo\gram}=\SI{1.982}{\kilo\gram}\).
\end{loigiai}

\begin{vidu}
Dung dịch \ce{HCl} có \(C_M=\SI{0.50}{\mole\per\liter}\). Trong \SI{250}{\milli\liter} dung dịch có bao nhiêu mol và bao nhiêu gam \ce{HCl}? Cho \(M_{\ce{HCl}}=36.5\) g/mol.
\end{vidu}
\begin{loigiai}
\[
\SI{250}{\milli\liter}=\SI{0.250}{\liter},\qquad
n=0.50\times0.250=\SI{0.125}{\mole}.
\]
\[
m=0.125\times36.5\approx\SI{4.56}{\gram}.
\]
\end{loigiai}

\subsection{Khối lượng riêng và bài toán dung dịch thực tế}
\begin{vidu}
Xăng E10 chứa \(10\%\) ethanol theo thể tích. Một xe có \SI{3.0}{\liter} E10. Tính thể tích và khối lượng ethanol, biết \(D_{\text{ethanol}}=\SI{0.789}{\gram\per\milli\liter}\).
\source{Dữ kiện theo bài nhiên liệu E10, đề chuyên Hóa Lâm Đồng 2026.}
\end{vidu}
\begin{loigiai}
\[
V_{\text{ethanol}}=3.0\times10\%=\SI{0.30}{\liter}=\SI{300}{\milli\liter}.
\]
\[
m=0.789\times300=\SI{236.7}{\gram}.
\]
\end{loigiai}

\begin{vidu}
Dung dịch \ce{H2SO4} \SI{98}{\percent} có khối lượng riêng \SI{1.84}{\gram\per\milli\liter}. Tính khối lượng \ce{H2SO4} nguyên chất trong \SI{1.00}{\liter} dung dịch.
\end{vidu}
\begin{loigiai}
\[
m_{\text{dd}}=1.84\times1000=\SI{1840}{\gram}.
\]
\[
m_{\ce{H2SO4}}=1840\times0.98=\SI{1803.2}{\gram}\approx\SI{1.80}{\kilo\gram}.
\]
\end{loigiai}

\subsection{Độ tan và kết tinh}
\begin{ghinho}
Độ tan \(S\) ở một nhiệt độ cho biết số gam chất tan tối đa trong \SI{100}{\gram} nước:
\[
S=\frac{m_{\text{chất tan}}}{m_{\text{nước}}}\times100.
\]
Khi làm lạnh dung dịch bão hòa, nếu độ tan giảm thì phần chất tan vượt quá độ tan mới sẽ tách ra.
\end{ghinho}

\begin{vidu}
Ở \SI{60}{\degreeCelsius}, độ tan của X là \(80\) g/\SI{100}{\gram} nước; ở \SI{20}{\degreeCelsius}, độ tan là \(30\) g/\SI{100}{\gram} nước. Dùng \SI{500}{\gram} nước pha dung dịch bão hòa ở \SI{60}{\degreeCelsius}, rồi làm lạnh đến \SI{20}{\degreeCelsius}. Giả sử nước không bay hơi, tính khối lượng X kết tinh.
\end{vidu}
\begin{loigiai}
\[
m_1=500\times\frac{80}{100}=\SI{400}{\gram},\qquad
m_2=500\times\frac{30}{100}=\SI{150}{\gram}.
\]
\[
m_{\text{kết tinh}}=400-150=\SI{250}{\gram}.
\]
\end{loigiai}

\begin{vidu}
Một cơ sở thu được \SI{1.20}{\ton} glucose và dùng toàn bộ để pha dung dịch G-5, trong đó \SI{100}{\milli\liter} dung dịch chứa \SI{5}{\gram} glucose. Mỗi chai chứa \SI{500}{\milli\liter}. Tính số chai sản xuất được.
\source{Phỏng theo Câu 5.1a, đề chuyên Hóa Lâm Đồng 2026.}
\end{vidu}
\begin{loigiai}
Mỗi chai chứa \(5\times500/100=\SI{25}{\gram}\) glucose. Đổi \SI{1.20}{\ton}=\(1.20\times10^6\) g:
\[
N=\frac{1.20\times10^6}{25}=48000.
\]
\end{loigiai}

\begin{tuluyen}
Pha \SI{5.0}{\kilo\gram} dung dịch \ce{NaCl} \(2.0\%\). Tính khối lượng \ce{NaCl} và nước.
\dapso{\SI{100}{\gram} \ce{NaCl}; \SI{4.90}{\kilo\gram} nước.}
\end{tuluyen}

\begin{tuluyen}
Trong \SI{2.0}{\liter} dung dịch \ce{NaOH} \(0.25\) M có bao nhiêu mol và bao nhiêu gam \ce{NaOH}? Cho \(M=40\) g/mol.
\dapso{\SI{0.50}{\mole}; \SI{20}{\gram}.}
\end{tuluyen}

\begin{tuluyen}
Một dung dịch có \(D=\SI{1.25}{\gram\per\milli\liter}\) và chứa \(20\%\) chất tan theo khối lượng. Tính khối lượng chất tan trong \SI{4.0}{\liter} dung dịch.
\dapso{\SI{1.00}{\kilo\gram}.}
\end{tuluyen}

\begin{tuluyen}
Độ tan của Y giảm từ \(60\) g xuống \(25\) g trên \SI{100}{\gram} nước khi làm lạnh. Dung dịch ban đầu chứa \SI{300}{\gram} nước và bão hòa ở nhiệt độ cao. Tính khối lượng Y kết tinh, giả sử không mất nước.
\dapso{\SI{105}{\gram}.}
\end{tuluyen}
